\section{Menguji independensi dalam tabel dua arah}
\label{twoWayTablesAndChiSquare}

\index{data!iPod|(}

\newcommand{\iPodAA}{2}
\newcommand{\iPodAB}{23}
\newcommand{\iPodAC}{36}
\newcommand{\iPodAD}{61}
\newcommand{\iPodAFraction}{0.2785}
\newcommand{\iPodAExpected}{20.33}
\newcommand{\iPodBA}{71}
\newcommand{\iPodBB}{50}
\newcommand{\iPodBC}{37}
\newcommand{\iPodBD}{158}
\newcommand{\iPodBFraction}{0.7215}
\newcommand{\iPodBExpected}{52.67}
\newcommand{\iPodDA}{73}
\newcommand{\iPodDB}{73}
\newcommand{\iPodDC}{73}
\newcommand{\iPodDD}{219}
\newcommand{\iPodN}{\iPodDD}

Kita semua pernah membeli produk bekas --
mobil, komputer, buku teks, dan sebagainya --
dan terkadang kita beranggapan bahwa penjual produk tersebut
akan berterus terang mengenai segala masalah tersembunyi pada
barang yang mereka jual.
Anggapan ini tidak seharusnya kita terima begitu saja.
Para peneliti merekrut \iPodN{} peserta dalam sebuah studi; mereka
diminta menjual iPod bekas\footnote{Bagi pembaca yang tidak setua
  para penulis, iPod pada dasarnya merupakan iPhone tanpa
  layanan seluler apa pun, dengan asumsi perangkat tersebut berasal dari salah satu
  generasi yang lebih akhir. Generasi-generasi awal lebih sederhana.}
yang diketahui pernah macet dua kali sebelumnya.
Para peserta diberi insentif untuk memperoleh uang sebanyak mungkin
dari iPod tersebut karena mereka akan menerima bagian sebesar 5\%
dari hasil penjualan, ditambah \$10 sebagai imbalan atas partisipasi.
Para peneliti ingin memahami jenis pertanyaan apa
yang akan mendorong penjual mengungkapkan masalah macet tersebut.

Tanpa sepengetahuan peserta yang menjadi penjual
dalam studi tersebut,
para pembeli bekerja sama dengan para peneliti
untuk menilai pengaruh berbagai pertanyaan
terhadap kemungkinan penjual mengungkapkan
masalah yang pernah terjadi pada iPod.
Pembeli yang mengikuti naskah memulai dengan
``Baiklah, sepertinya saya yang harus mulai dulu.
  Jadi, Anda sudah memiliki iPod ini selama 2 tahun ...''
dan mengakhirinya dengan salah satu dari tiga pertanyaan:
\begin{itemize}
\item Umum: Apa yang bisa Anda ceritakan tentang perangkat ini?
\item Asumsi positif: Perangkat ini tidak bermasalah, bukan?
\item Asumsi negatif: Masalah apa yang dimiliki perangkat ini?
\end{itemize}
Pertanyaan tersebut merupakan perlakuan yang diberikan kepada para penjual,
sedangkan responsnya adalah apakah pertanyaan itu mendorong mereka
untuk mengungkapkan masalah macet pada iPod.
Hasilnya ditampilkan pada Gambar~\ref{ipod_ask_data_summary},
dan data tersebut menunjukkan bahwa pertanyaan
\emph{Masalah apa yang dimiliki perangkat ini?}
merupakan yang paling efektif dalam mendorong penjual mengungkapkan
masalah macet yang pernah terjadi.
Namun, Anda juga perlu bertanya kepada diri sendiri:
mungkinkah hasil ini muncul semata-mata karena kebetulan,
atau apakah hasil ini benar-benar menjadi bukti bahwa beberapa pertanyaan
lebih efektif untuk mengungkap kebenaran?

\begin{figure}[ht]
\centering
\begin{tabular}{l ccc l}
  \hline
  & Umum & Asumsi positif &
      Asumsi negatif & Total \\ 
  \hline
  Mengungkap masalah & \iPodAA{} &  \iPodAB{} &
      \iPodAC{} & \iPodAD{} \\ 
  Menyembunyikan masalah &  \iPodBA{} &  \iPodBB{} &
      \iPodBC{} & \iPodBD{} \\ 
  \hline
  Total & \iPodDA{} & \iPodDB{} &
      \iPodDC{} & \iPodDD{} \\
  \hline
\end{tabular}
\caption{Ringkasan studi iPod, dengan sebuah pertanyaan
  diajukan kepada peserta studi yang bertindak sebagai penjual.}
\label{ipod_ask_data_summary}
\end{figure}

\begin{onebox}{Perbedaan antara tabel satu arah dan tabel dua arah}
  Tabel satu arah menggambarkan cacah untuk setiap hasil dalam satu
  variabel.
  Tabel dua arah menggambarkan cacah untuk \emph{kombinasi}
  hasil dari dua variabel.
  Ketika meninjau tabel dua arah, kita sering ingin mengetahui,
  apakah variabel-variabel tersebut saling berkaitan?
  Dengan kata lain, apakah keduanya dependen (alih-alih independen)?
\end{onebox}

Uji hipotesis untuk eksperimen iPod sebenarnya bertujuan
menilai apakah terdapat bukti yang signifikan secara statistik
bahwa tingkat keberhasilan setiap pertanyaan dalam mendorong peserta
untuk mengungkapkan masalah iPod berbeda-beda.
Dengan kata lain, tujuannya adalah memeriksa apakah jenis
pertanyaan pembeli independen dari pengungkapan
masalah oleh penjual.


\D{\newpage}

\subsection{Cacah harapan dalam tabel dua arah}

\noindent%
Seperti pada tabel satu arah, kita perlu menghitung
cacah harapan untuk setiap sel dalam tabel dua arah.

\begin{examplewrap}
\begin{nexample}{Dari eksperimen tersebut,
    kita dapat menghitung proporsi seluruh penjual yang mengungkapkan
    masalah macet sebagai $\iPodAD{}/\iPodDD = \iPodAFraction{}$.
    Jika benar-benar tidak ada perbedaan di antara pertanyaan
    dan 27.85\% penjual memang akan mengungkapkan masalah
    macet apa pun pertanyaan yang diajukan kepada mereka,
    berapa banyak dari \iPodDA{} orang dalam kelompok \resp{Umum}
    yang kita harapkan akan mengungkapkan masalah
    macet tersebut?} \label{iPodExComputeExpAA}
  Kita memperkirakan bahwa $\iPodAFraction{} \times \iPodDA{} = \iPodAExpected{}$
  penjual akan mengungkapkan masalah tersebut.
  Jelas bahwa cacah teramati lebih kecil dari nilai ini, meskipun belum
  jelas apakah hal tersebut disebabkan oleh variasi acak atau karena
  pertanyaan-pertanyaan itu berbeda dalam hal seberapa efektif
  pertanyaan tersebut mengungkap kebenaran.
\end{nexample}
\end{examplewrap}

\begin{exercisewrap}
\begin{nexercise}\label{iPodExComputeExpBB}
Jika pertanyaan-pertanyaan tersebut sebenarnya sama efektifnya,
yakni sekitar 27.85\% responden akan mengungkapkan
masalah macet apa pun pertanyaan yang diajukan kepada mereka,
kira-kira berapa penjual dari kelompok Asumsi positif
yang kita harapkan akan \emph{menyembunyikan} masalah
macet tersebut?\footnotemark
\end{nexercise}
\end{exercisewrap}
\footnotetext{Kita mengharapkan
    $(1 - \iPodAFraction{}) \times \iPodDA{} = \iPodBExpected{}$.
    Tidak masalah jika hasil ini,
    seperti hasil dari Contoh~\ref{iPodExComputeExpAA},
    berupa pecahan.}

Kita dapat menghitung banyaknya penjual yang kita harapkan
akan mengungkapkan atau menyembunyikan masalah macet untuk semua kelompok,
jika pertanyaan tersebut tidak memengaruhi apa yang mereka ungkapkan,
dengan menggunakan strategi yang sama seperti pada
Contoh~\ref{iPodExComputeExpAA} dan
Latihan Terbimbing~\ref{iPodExComputeExpBB}.
Cacah harapan ini digunakan untuk menyusun Gambar~\ref{ipod_ask_data_summary_expected},
yang sama dengan Gambar~\ref{ipod_ask_data_summary},
hanya saja kini cacah harapan telah ditambahkan di dalam tanda kurung.

\begin{figure}[h]
\centering
\begin{tabular}{l lll l}
  \hline
  & Umum & Asumsi positif &
      Asumsi negatif & Total \\ 
  \hline
  Mengungkap masalah &
      \iPodAA{} \ \highlightO{\footnotesize(\iPodAExpected{})} &
      \iPodAB{} \highlightO{\footnotesize(\iPodAExpected{})} &
      \iPodAC{} \highlightO{\footnotesize(\iPodAExpected{})} &
      \iPodAD{} \\ 
  Menyembunyikan masalah &
      \iPodBA{} \highlightO{\footnotesize(\iPodBExpected{})} &
      \iPodBB{} \highlightO{\footnotesize(\iPodBExpected{})} &
      \iPodBC{} \highlightO{\footnotesize(\iPodBExpected{})} &
      \iPodBD{} \\ 
  \hline
  Total & \iPodDA{} & \iPodDB{} &
      \iPodDC{} & \iPodDD{} \\
  \hline
\end{tabular}
\caption{Cacah teramati dan
    \highlightO{(cacah harapan)}.}
\label{ipod_ask_data_summary_expected}
\end{figure}

Contoh dan latihan di atas memberikan beberapa petunjuk
untuk menghitung cacah harapan.
Secara umum, cacah harapan untuk tabel dua arah dapat
dihitung menggunakan total baris, total kolom,
dan total tabel.
Misalnya, jika tidak ada perbedaan di antara kelompok,
maka sekitar 27.85\% dari setiap kolom seharusnya berada pada baris pertama:
\begin{align*}
\iPodAFraction{}\times (\text{total kolom 1}) &= \iPodAExpected{} \\
\iPodAFraction{}\times (\text{total kolom 2}) &= \iPodAExpected{} \\
\iPodAFraction{}\times (\text{total kolom 3}) &= \iPodAExpected{}
\end{align*}
Dengan mengingat kembali cara menghitung \iPodAFraction{} --
sebagai fraksi penjual yang mengungkapkan masalah macet
($\iPodAD{}/\iPodDD{}$) --
ketiga cacah harapan ini dapat dihitung sebagai
\begin{align*}
\left(\frac{\text{total baris 1}}{\text{total tabel}}\right)
    \text{(total kolom 1)} &= \iPodAExpected{} \\
\left(\frac{\text{total baris 1}}{\text{total tabel}}\right)
    \text{(total kolom 2)} &= \iPodAExpected{} \\
\left(\frac{\text{total baris 1}}{\text{total tabel}}\right)
    \text{(total kolom 3)} &= \iPodAExpected{}
\end{align*}
Hal ini menghasilkan rumus umum untuk menghitung cacah
harapan dalam tabel dua arah ketika kita ingin menguji apakah
terdapat bukti kuat mengenai hubungan antara variabel
kolom dan variabel baris.

\D{\newpage}

\begin{onebox}{Menghitung cacah harapan dalam tabel dua arah}
  Untuk menentukan cacah harapan bagi baris ke-$i$
  dan kolom ke-$j$, hitung
  \begin{align*}
  \text{cacah harapan}_{\text{baris }i,\text{ kolom }j}
    = \frac{(\text{total baris $i$}) \times
        (\text{total kolom $j$})}{\text{total tabel}}\vspace{2mm}
  \end{align*}
\end{onebox}

%
